\documentclass[10pt,a4paper]{amsart}
\usepackage[margin=19mm]{geometry}
\usepackage{amsmath,amssymb,mathtools}
\usepackage[scr=rsfs,cal=euler]{mathalfa}
\usepackage{xfrac}
\usepackage[unicode,hidelinks,bookmarks=false]{hyperref}
\newcommand{\ie}{i.e.\ }
\newcommand{\cf}{cf.\ }
\newcommand{\resp}{respectively\ }
\newcommand{\Subsection}{\subsection}
\providecommand{\nobreakdash}{\nobreak\hbox{-}}
\setlength{\emergencystretch}{2em}
\multlinegap0pt
\usepackage{bm}
\usepackage{bbm}
\usepackage{dsfont}
\usepackage{xspace}
\usepackage{cancel}
\usepackage{mathtools}
\usepackage{amsthm}
\makeatletter
\@ifundefined{Theorem}{\newtheorem{Theorem}{Theorem}}{}
\@ifundefined{Lemma}{\newtheorem{Lemma}[Theorem]{Lemma}}{}
\@ifundefined{Remark}{\newtheorem{Remark}[Theorem]{Remark}}{}
\makeatother
\title[Gelfand-Shilov spaces for extended Gevrey regularity]{Gelfand-Shilov spaces for extended Gevrey regularity}
\author{Nenad Teofanov, Filip Tomic, Milica Zigic}
\date{Source: arXiv:2503.13634v1 (CC BY 4.0)}
\begin{document}
\maketitle

\noindent\emph{Source and licence.} Gelfand-Shilov spaces for extended Gevrey regularity. Version arXiv:2503.13634v1 (CC BY 4.0), distributed under the Creative Commons Attribution 4.0 International licence, as recorded in the arXiv licence metadata for this exact version and shown on the abstract page https://arxiv.org/abs/2503.13634v1. Full rights notice: licenses/arxiv-2503.13634-CC-BY-4.0.txt.\par\smallskip

\noindent\textit{Editorial scope.} Counted unit: the Theorem whose source label is \texttt{nec-suf-cond} in the article (source0.tex lines 1296--1305, source label \texttt{nec-suf-cond}), delivered together with the article's own complete original proof of it (source0.tex lines 1307--1501, one \texttt{proof} environment of 195 lines). No mathematics is written, repaired or re-derived: statement and proof are the article's own text line for line. Required context retained uncounted: eq:the-weight-matrix (lines 170--172); M.1' (lines 220--222); thm:with-or-without-h (lines 480--494); rem:alfa-beta (lines 550--555); thm:L2-and-Linfty (lines 589--598); eq:Linfty-implies-L2 (lines 680--701); Thm:opisFurije (lines 1037--1048); GRandRelatives (lines 1233--1242); eq:inversion (lines 1262--1265). Every definition, display and label of those lines is retained unchanged. Omitted from this package and not redistributed: 1--169, 173--219, 223--479, 495--549, 556--588, 599--679, 702--1036, 1049--1232, 1243--1261, 1266--1295, 1306--1306, 1502--1753. Nothing inside a retained excerpt is omitted or altered: every retained body is delivered exactly as the article's lines are written, so no omission receipt of any kind accompanies this package. Citations: the retained excerpts carry no citation command at all, so no bibliography entry of the article is transcribed and no reference is redistributed. Reference adaptation: none was needed and none was made; every reference inside the delivered unit resolves inside it, so no unresolved reference or citation is produced. Printed slips kept literal: no printed word, symbol, hypothesis or number of the counted statement or of its proof was altered. Layout and class: the article's own class is \texttt{article} with packages including graphicx, multirow, amsmath, amssymb, amsfonts, amsthm, mathrsfs, appendix, xcolor, textcomp, manyfoot, booktabs, algorithm, algorithmicx; the wrapper is the lane's pinned \texttt{amsart} editorial wrapper at 10pt on A4, which restates the theorem-like environments the retained text needs and adds the article's own macro definitions that the retained text needs (none), with extra wrapper packages (bm, bbm, dsfont, xspace, cancel, mathtools). The delivered numbering is the unit's own. Assets: no retained span contains a figure environment, a graphics-inclusion command, a PDF-page inclusion, a table or a TikZ picture; no image file is copied into this package, the package declares no asset dependency and no image file is redistributed with it. Licence: version arXiv:2503.13634v1 of this article is distributed under the Creative Commons Attribution 4.0 International licence, as recorded in the arXiv licence metadata for this exact version and shown on the abstract page https://arxiv.org/abs/2503.13634v1. A full rights notice is delivered at licenses/arxiv-2503.13634-CC-BY-4.0.txt. Characters: every retained excerpt is pure ASCII, so the delivered engine, XeLaTeX with its default Latin Modern text font, reports no missing character.
    \begin{equation}\label{eq:the-weight-matrix}
\mathcal{M}_\sigma=\left\{M^{(\tau)} =  (M_{p}^{\tau,\sigma})_{p\in\mathbb N_0}\colon \tau>0\right\}.
    \end{equation}

\begin{equation} \label{M.1'}
M_p^{\tau,\sigma}M_q^{\tau,\sigma}\leq M_{p+q}^{\tau,\sigma},\qquad p,q\in\mathbb N_0,
\end{equation}

\begin{Theorem} \label{thm:with-or-without-h}
Let $\sigma>1$, and let the weight matrix $\mathcal{M}_\sigma$ be given by 
\eqref{eq:the-weight-matrix}, 
$ M_{p}^{\tau,\sigma}=p^{\tau p^\sigma}$, $ p\in \mathbb{N}$, $ M_0^{\tau,\sigma}=1$.
Then 
\begin{equation}
\mathcal S_{[\mathcal{M}_\sigma ] }^{ [\mathcal{M}_\sigma]} (\mathbb R^n)
= \mathscr S_{[\mathcal{M}_\sigma ] }^{ [\mathcal{M}_\sigma ]} (\mathbb R^n),    
\end{equation}
$\mathcal S^{ [\mathcal{M}_\sigma]} (\mathbb R^n)
= \mathscr S^{ [\mathcal{M}_\sigma ]} (\mathbb R^n),$
and 
$\mathcal S_{[\mathcal{M}_\sigma ] } (\mathbb R^n)
= \mathscr S_{[\mathcal{M}_\sigma ] } (\mathbb R^n).$
\end{Theorem}

\begin{Remark} \label{rem:alfa-beta}
By  $\widetilde{(M.2)}$, \eqref{M.1'}, and Theorem \ref{thm:with-or-without-h}
it follows that the product $M^{\tau, \sigma} _{|\alpha|} 
  M^{\tau, \sigma} _{|\beta|}$ in the definition of 
$\mathcal S_{[\mathcal{M}_\sigma ] }^{ [\mathcal{M}_\sigma]} (\mathbb R^n)$ can be replaced by $M^{\tau, \sigma} _{|\alpha + \beta|}$ yielding the same spaces.
 \end{Remark}

\begin{Theorem} \label{thm:L2-and-Linfty}
Let $\sigma>1$, and let the weight matrix $\mathcal{M}_\sigma$ be given by 
\eqref{eq:the-weight-matrix}, 
$ M_{p}^{\tau,\sigma}=p^{\tau p^\sigma}$, $ p\in \mathbb{N}$, $ M_0^{\tau,\sigma}=1$.
Then 
\begin{equation}
\mathcal S_{[\mathcal{M}_\sigma ] }^{ [\mathcal{M}_\sigma]} (\mathbb R^n)
=  \mathcal S_{[\mathcal{M}_\sigma],2}^{[\mathcal{M}_\sigma]} (\mathbb R^n).    \label{eq:L2-and-Linfty}
\end{equation}
\end{Theorem}

\begin{align}
\| x^{\alpha } \partial ^{\beta }  \varphi(x)\|_{L^2} 
& \leq C_s  \sup_{x} (1+|x|^2)^{s} |x^\alpha \partial^\beta \varphi(x)| \nonumber \\
& = C_s  \sum_{|\gamma| \leq s} 
C_\gamma  \sup_{x} | x^{\alpha + 2\gamma} \partial^\beta \varphi(x)| 
\nonumber \\
%& \leq C_1  \sum_{|\gamma| \leq s}  \sup_{x} (1+|x|^2)^{s} |x^\alpha 
%\partial^\beta \varphi(x)| \nonumber \\
& \leq C_1  \sum_{|\gamma| \leq s} 
 M^{\tau,\sigma}_{|\alpha + 2\gamma|}   M^{\tau,\sigma}_{|\beta|} \nonumber \\
& \leq C_1  \sum_{|\gamma| \leq s} C_2 ^{(2\gamma) 2^{\sigma - 1}(|\alpha |^\sigma + (2\gamma) ^\sigma)}
M^{\tau,\sigma}_{|\alpha |}   M^{\tau,\sigma}_{|\beta|} 
\nonumber \\
& \leq C_1  C_2 ^{ 2^{\sigma - 1}(2s) ^{\sigma+1}}
\sum_{|\gamma| \leq s} h^{|\alpha |^\sigma }
M^{\tau,\sigma}_{|\alpha |}   M^{\tau,\sigma}_{|\beta|}  
\nonumber \\
& \leq C
h^{|\alpha |^\sigma }
M^{\tau,\sigma}_{|\alpha |}   M^{\tau,\sigma}_{|\beta|},
\label{eq:Linfty-implies-L2}
\end{align}

\begin{Theorem} \label{Thm:opisFurije}
Let $\sigma>1$,  and let the weight matrix $\mathcal{M}_\sigma$ be given by \eqref{eq:the-weight-matrix}, 
$ M_{p}^{\tau,\sigma}=p^{\tau p^\sigma}$, $ p\in \mathbb{N}$, $ M_0^{\tau,\sigma}=1$.
Then the following conditions are equivalent:
    \begin{enumerate}
        \item $\varphi\in \mathcal  S_{\{\mathcal{M}_\sigma\}}^{\{\mathcal{M}_\sigma\}} (\mathbb R^n);$
        \item $(\exists \tau>0)(\exists C>0)\;(\forall \alpha,\beta\in \mathbb N_0^n)$
        $$\sup_x|x^\alpha \varphi(x)|\leq C M^{\tau,\sigma}_{|\alpha|} \quad \text{and}\quad \sup_x|\partial^\beta \varphi(x)|\leq CM^{\tau,\sigma}_{|\beta|};$$
        \item $(\exists \tau>0)(\exists C>0)\;(\forall \alpha,\beta\in \mathbb N_0^n)$
        $$\sup_x|x^\alpha \varphi(x)|\leq C M^{\tau,\sigma}_{|\alpha|} \quad \text{and}\quad \sup_\xi|\xi^\beta \widehat\varphi(\xi)|\leq CM^{\tau,\sigma}_{|\beta|}.$$
    \end{enumerate}
\end{Theorem}

\begin{Lemma} \label{GRandRelatives}
Let  $f, g \in L^2 (\mathbb{R}^n).$ Then we have:
\begin{eqnarray*}
W(f,g) (x,\omega)  & = & 2^n R_g f (x,\omega), \\
V_g f (x,\omega) & = &  e^{-\pi i x \omega} R_{\check{g}} f \left(\frac{x}{2},\frac{\omega}{2}\right),
\\
A(f,g) (x,\omega) & = & R_{\check{g}} f \left(\frac{x}{2},\frac{\omega}{2}\right), \quad
x,\omega \in \mathbb{R}^n.
\end{eqnarray*}
\end{Lemma}

\begin{equation} \label{eq:inversion}
f = \frac{1}{\langle g_2, g_1 \rangle} \iint  R_{g_1} f (x, \omega)
R g_2 (x,\omega) \, dx d\omega.
\end{equation}

\begin{Theorem} \label{nec-suf-cond}
Let $\sigma>1$, and let the weight matrix $\mathcal{M}_\sigma$ be given by 
\eqref{eq:the-weight-matrix}, 
$ M_{p}^{\tau,\sigma}=p^{\tau p^\sigma}$, $ p\in \mathbb{N}$, $ M_0^{\tau,\sigma}=1$.
If $f,g \in    \mathcal S_{[\mathcal{M}_\sigma ]}^{[\mathcal{M}_\sigma]} (\mathbb R^{n}) $, 
and  $ T \in TFA$, then $ T   \in   \mathcal S_{[\mathcal{M}_\sigma ]}^{[\mathcal{M}_\sigma]} (\mathbb R^{2n})$.

Conversely, if  $ T \in TFA$ and $ T   \in   \mathcal S_{[\mathcal{M}_\sigma ]}^{[\mathcal{M}_\sigma]} (\mathbb R^{2n})$
then $f,g \in     \mathcal S_{[\mathcal{M}_\sigma ]}^{[\mathcal{M}_\sigma]} (\mathbb R^{n}).$
\end{Theorem}

\begin{proof}
Since $\mathcal S_{[\mathcal{M}_\sigma ]}^{[\mathcal{M}_\sigma]} (\mathbb R^{n})$
are closed under reflections, dilations and modulations, then by Lemma \ref{GRandRelatives}
it suffices to give the proof for any $T \in TFA$.
We will give the proof for the Grossmann-Royer transform, and the same conclusion holds for other time-frequency representations in question. 

\par

Put $ \Phi (x,t) := f(2x -t) g (t)$, $x,t \in \mathbb R^{n}$.
Then 
$$
R_g f (x,\omega) = e^{-4\pi i \omega x}  \int e^{2\pi i \omega (2t) } \Phi (x,t)\, dt. 
$$
The first part of the claim will be proved 
if we show that
$ \Phi \in \mathcal S_{[\mathcal{M}_\sigma ]}^{[\mathcal{M}_\sigma]} (\mathbb R^{2n}) $,
since $  \mathcal S_{[\mathcal{M}_\sigma ]}^{[\mathcal{M}_\sigma]} (\mathbb R^{2n})$ is closed under dilations and modulations,
and since then $\mathcal F_2 \Phi \in 
  \mathcal S_{[\mathcal{M}_\sigma ]}^{[\mathcal{M}_\sigma]} (\mathbb R^{2n})$.
  
As before, we give the proof for the Beurling case, and the 
Roumieu case follows by similar arguments.

By Theorem \ref{Thm:opisFurije} it is enough to show
\begin{equation} \label{varphi po x}
\sup_{x,t \in \mathbb{R}^n}  
| x^{\alpha}  t^\beta  \Phi (x,t) | \leq C 
M^{\tau, \sigma} _{|\alpha|} M^{\tau, \sigma} _{|\beta|},\quad \alpha,\beta\in\mathbb N_0^n,
\end{equation}
and
\begin{equation} \label{varphi po ksi}
\sup_{x,t \in \mathbb{R}^n}  
|  \partial^\alpha _x \partial^\beta _t  \Phi (x,t) |
\leq C M^{\tau, \sigma} _{|\alpha|} M^{\tau, \sigma} _{|\beta|},\quad \alpha,\beta\in\mathbb N_0^n,
\end{equation}
for any given $\tau > 0$, and some constant $C= C_\tau >0$.

The decay of  $ \Phi (x,t) := f(2x -t) g (t)$ can be estimated as follows.
%%%%%%%
\begin{align}
\sup_{x,t \in \mathbb{R}^d}  | x^{\alpha}  t^\beta  f(2x-t) g (t) |
& \leq  
2^{-|\alpha|}\sup_{y,t \in \mathbb{R}^n}  |y+t|^{|\alpha|}  |t|^{|\beta|} 
|f(y) | |  g (t) | 
\nonumber \\
&\leq C 
2^{-|\alpha|}\sup_{y,t \in \mathbb{R}^n }
\sum_{\gamma \leq \alpha} \binom{\alpha}{\gamma}
|y|^{|\alpha|}  |t|^{|\alpha - \gamma| + |\beta|} 
|f(y) | |  g (t) | 
\nonumber \\
&\leq C  
\sup_{y \in \mathbb{R}^n} |y|^{|\alpha|}  |f(y) | \cdot
\sup_{t \in \mathbb{R}^n} |t|^{|\alpha| + |\beta|} 
 |  g (t) | 
\nonumber \\
& \leq
C  M^{\tau/2,\sigma}_{|\alpha |}  
M^{\tau/2,\sigma}_{|\alpha | +|\beta|} 
%\leq C  M^{\tau/2,\sigma}_{|\alpha |}  M^{\tau/2,\sigma}_{|\alpha | +|\beta|} 
\nonumber \\
& \leq
C  M^{\tau,\sigma}_{|\alpha | + |\beta|}, 
\label{eq:est-mult}
\end{align}
see also Remark \ref{rem:alfa-beta}.

%%%%%%


For the proof of  \eqref{varphi po ksi}, we use the Leibniz formula
and properties of the sequence $(M_{p}^{\tau,\sigma})_{p\in\mathbb N_0}$:

\begin{eqnarray*}
 \sup_{x,t \in \mathbb{R}^n} |\partial^\alpha _x \partial^\beta _t   f(2x-t) g (t)|
& \leq &
 2^{|\alpha|} \sum_{\gamma \leq \beta}
  \binom{\beta}{\gamma} 
\sup_{x,t \in \mathbb{R}^n} |\partial^{\alpha} _x \partial^{\gamma} _t  f (2x - t)
 \partial^{\beta - \gamma} _t g  (t)|
\\
& \leq & 2^{|\alpha| }  \sum_{\gamma \leq \beta}
  \binom{\beta}{\gamma} 
\sup_{x,t \in \mathbb{R}^n}
| \partial^{\alpha} _x \partial^{\gamma} _t  f (2x - t)|
\sup_{t \in \mathbb{R}^n} |\partial^{\beta - \gamma} _t g  (t)|
\\
& \leq & 2^{|\alpha| +|\beta|} 
C M^{\tau/4,\sigma}_{|\alpha |  +|\beta|}  
M^{\tau/4,\sigma}_{|\beta|} 
\leq  C_1 M^{\tau,\sigma}_{|\alpha |  +|\beta|}, 
\end{eqnarray*}
for some constants $C, C_1 > 0 $.

Thus $ R_g f \in \mathcal S_{(\mathcal{M}_\sigma )}^{(\mathcal{M}_\sigma)} (\mathbb R^{2n}) $ if  $f,g \in    \mathcal S_{(\mathcal{M}_\sigma )}^{(\mathcal{M}_\sigma)} (\mathbb R^{n}) .$

%Therefore, $ \varphi \in \in \mathcal S_{[\mathcal{M}_\sigma ]}%^{[\mathcal{M}_\sigma]} (\mathbb R^{2n}) $ and, consequently,
%$ R_g f (x,\o) \in  
%\mathcal S_{[\mathcal{M}_\sigma ]}^{[\mathcal{M}_\sigma]} (\mathbb %R^{2n}) $.

Now, assume that 
\begin{equation} \label{eq:RinGSspace}
R_g f \in  
\mathcal S_{(\mathcal{M}_\sigma )}^{(\mathcal{M}_\sigma)} (\mathbb R^{2n}) 
\end{equation}
for a given $g \in L^2 (\mathbb R^{2n}) \setminus \{ 0\}$.

By the inversion formula \eqref{eq:inversion}
we have
\begin{equation} \label{eq:term-for-estimate}
t^\alpha \partial^\beta _t f (t) = \frac{1}{\langle h, g \rangle} 
\iint R_{g} f (x, \omega)
t^\alpha \partial^\beta _t (R (h (t)) (x, \omega)\, dx d\omega,
\end{equation}
where we may choose $h \in \mathcal S_{(\mathcal{M}_\sigma )}^{(\mathcal{M}_\sigma)} (\mathbb R^{n})$ such that 
$\langle h, g \rangle = 1$.

Let there be given $\tau >0$.

From \eqref{eq:RinGSspace} we get
\begin{eqnarray}
% \sup_{t \in \mathbb{R}^n} 
 | t^\alpha \partial^\beta _t (R (h (t)) (x, \omega) |
&  = &
%  \sup_{t \in \mathbb{R}^n} 
| t^\alpha \partial^\beta _t  ( e^{4\pi i \omega (t-x)} h(2x-t))|
\nonumber \\
&  \leq &
C \sum_{\gamma \leq \beta} \binom{\beta}{\gamma}
|\omega^{\beta - \gamma} |
 | t^{\alpha}  \partial^\gamma _t h(2x-t))| 
 \nonumber \\
&  \leq &
C M^{\tau/2, \sigma} _{|\alpha|} M^{\tau/2, \sigma} _{|\beta|} 
\sum_{\gamma \leq \beta} \binom{\beta}{\gamma}
|\omega^{\beta - \gamma} |, 
\nonumber
%\label{eq:Restimate-1}
\end{eqnarray}
wherefrom
\begin{equation} \label{eq:term-for-estimate-2}
|t^\alpha \partial^\beta _t f (t) |
\leq C
M^{\tau/2, \sigma} _{|\alpha|} M^{\tau/2, \sigma} _{|\beta|} 
\sum_{\gamma \leq \beta} \binom{\beta}{\gamma}
\iint |\omega^{\beta - \gamma}  R_{g} f (x, \omega)|
dx d\omega.
\end{equation}
Next we estimate the double integral in \eqref{eq:term-for-estimate-2}. We use  similar argument as in the proof of Theorem \ref{thm:L2-and-Linfty} (cf. \eqref{eq:Linfty-implies-L2}):
\begin{align}
\iint |\omega^{\beta - \gamma}  R_{g} f (x, \omega) & |
 dx d\omega \nonumber \\
  & \leq 
C_1 \sup_{x,\omega \in \mathbb R^{n}} (1+|x|)^2
(1+|\omega|)^{2+|\beta - \gamma|} |R_{g} f (x, \omega)|
\nonumber \\
  & \leq 
C_2    M^{\tau/4, \sigma} _{2} M^{\tau/4, \sigma} _{2+|\beta -\gamma|}  \leq C_3 M^{\tau/4, \sigma} _{2+|\beta|} 
\nonumber \\
  & \leq 
C_4 ^{(2^{\sigma - 1} +1) |\beta|^\sigma} 
 M^{\tau/4, \sigma} _{|\beta|} 
 \nonumber
 % \label{eq:Restimate-2},
\end{align}
for some constants $ C_j >0,$ $ j =1,2,3,4$, wherefrom
\begin{equation}
\sum_{\gamma \leq \beta} \binom{\beta}{\gamma}
\iint |\omega^{\beta - \gamma}  R_{g} f (x, \omega)|
dx d\omega
\leq 
C_5 ^{(2^{\sigma - 1} +1) |\beta|^\sigma} 
M^{\tau/4, \sigma} _{|\beta|} 
\leq
 C M^{\tau/2, \sigma} _{|\beta|},
  \label{eq:Restimate-2}
\end{equation}
for some constants $C, C_5 > 0$.
By \eqref{eq:term-for-estimate-2} and 
\eqref{eq:Restimate-2} we get
$$
\sup_{t \in \mathbb R^{n}}
|t^\alpha \partial^\beta _t f (t) | \leq
C 
M^{\tau/2, \sigma} _{|\alpha|} M^{\tau/2, \sigma} _{|\beta|}  M^{\tau/2, \sigma} _{|\beta|}
\leq
C 
M^{\tau, \sigma} _{|\alpha|} M^{\tau, \sigma} _{|\beta|},  
$$
i.e. $ f \in \mathcal S_{(\mathcal{M}_\sigma )}^{(\mathcal{M}_\sigma)} (\mathbb R^{n}) $.

Finally, $ g \in \mathcal S_{(\mathcal{M}_\sigma )}^{(\mathcal{M}_\sigma)} (\mathbb R^{n}) $
follows from 
$  R_g f = \overline{R_f g}$.
\end{proof}

\end{document}
